\section{Clipping}
\label{sec:clip}
\here{5}
\concept{Primitives outside the view volume must be removed and primitives crossing it cut, before rasterisation.
The GPU clips every triangle against the six planes $-w_c\le x_c,y_c,z_c\le w_c$. The program adds two CPU steps:
whole parts are culled against the frustum before they are submitted, and grass is culled per cell.}
\begin{equation}
\text{plane } i:\ \vect a_i\cdot\vect p+d_i\ge0\ \text{inside};\qquad
\text{sphere }(\vect c,r)\ \text{rejected if}\ \exists i:\ \vect a_i\cdot\vect c+d_i<-r
\label{eq:cull}
\end{equation}
The six planes come from the rows of $PV$ (Gribb--Hartmann): left $=\text{row}_4+\text{row}_1$, right $=\text{row}_4-\text{row}_1$, and so on.

\begin{figure}[H]
\centering
\begin{subfigure}{0.49\linewidth}\raw{scene/v10_proj_front}{\linewidth}\caption{Near/far clipping in an orthographic front view.}\end{subfigure}\hfill
\begin{subfigure}{0.49\linewidth}\raw{lab/explorer_stage_03_assembly}{\linewidth}\caption{Frame Explorer: primitive assembly and the frustum.}\end{subfigure}
\caption{Clipping in the program. In (a) the view volume is an 80\,m slab around the camp: notice that the
backdrop hills and the far terrain are cut flat at the near and far planes, and the cut faces show the inside of
the terrain mesh, exactly as clipping produces (no capping). (b) shows the assembled triangles of the frame.}
\label{fig:clip}
\end{figure}

\begin{figure}[H]
\centering\wb{Clipping}{0.95\linewidth}
\caption{Workbook sheet \code{Clipping}: Cohen--Sutherland outcodes, Liang--Barsky parameters and
Sutherland--Hodgman steps. These are worked textbook examples computed by formulas, not code in the program.}
\label{fig:wb-clip}
\end{figure}

\begin{table}[H]
\centering\small
\begin{tabularx}{\linewidth}{@{}lXXl@{}}
\toprule
Method & What it does & Cost & Used in the program\\\midrule
Cohen--Sutherland & 4-bit outcodes, trivial accept/reject, then intersect & cheap for mostly inside/outside lines & no (workbook)\\
Liang--Barsky & parametric $t$ range, at most one division per edge & fewer intersections & no (workbook)\\
Sutherland--Hodgman & polygon against one plane at a time & output grows per plane & by the GPU (in clip space)\\
Frustum culling & whole object against 6 planes (Eq.~\ref{eq:cull}) & 6 dot products per part & yes, \code{math/Frustum}\\
\bottomrule
\end{tabularx}
\caption{Clipping methods compared.}
\end{table}
\where{src/math/Frustum.{h,cpp} (plane extraction, sphere test); grass cell culling in src/rendering/Grass.cpp}
\critique{There is no interactive user-defined clip window (the brief's item); clipping is shown through the
orthographic slab and the workbook. Frustum culling with bounding spheres is conservative: a long thin part (a
wall) has a large sphere and is drawn even when only its sphere touches the frustum.}
